\documentclass[tikz,border=8pt]{standalone}
\usepackage{ctex,amsmath,amssymb}
\usetikzlibrary{arrows.meta,positioning,calc,fit,backgrounds}
% Shared visual language for LFS architecture diagrams.
\RequirePackage{../../mymath}

\definecolor{LFSBlue}{HTML}{2563EB}
\definecolor{LFSBlueSoft}{HTML}{DBEAFE}
\definecolor{LFSGreen}{HTML}{15803D}
\definecolor{LFSGreenSoft}{HTML}{DCFCE7}
\definecolor{LFSOrange}{HTML}{C2410C}
\definecolor{LFSOrangeSoft}{HTML}{FFEDD5}
\definecolor{LFSRed}{HTML}{B91C1C}
\definecolor{LFSRedSoft}{HTML}{FEE2E2}
\definecolor{LFSPurple}{HTML}{7E22CE}
\definecolor{LFSPurpleSoft}{HTML}{F3E8FF}
\definecolor{LFSGray}{HTML}{475569}
\definecolor{LFSGraySoft}{HTML}{F1F5F9}

\tikzset{
  lfs picture/.style={
    font=\small,
    node distance=8mm and 12mm,
    >=Latex,
    every path/.style={line cap=round, line join=round}
  },
  block/.style={
    draw=LFSBlue,
    fill=LFSBlueSoft,
    rounded corners=2pt,
    line width=0.8pt,
    align=center,
    inner xsep=7pt,
    inner ysep=5pt,
    minimum height=10mm
  },
  compute/.style={
    block,
    draw=LFSPurple,
    fill=LFSPurpleSoft
  },
  storage/.style={
    block,
    draw=LFSGreen,
    fill=LFSGreenSoft
  },
  interface/.style={
    block,
    draw=LFSOrange,
    fill=LFSOrangeSoft
  },
  risk/.style={
    block,
    draw=LFSRed,
    fill=LFSRedSoft
  },
  neutral/.style={
    block,
    draw=LFSGray,
    fill=LFSGraySoft
  },
  decision/.style={
    diamond,
    draw=LFSOrange,
    fill=LFSOrangeSoft,
    line width=0.8pt,
    align=center,
    inner sep=2pt,
    aspect=2
  },
  group/.style={
    draw=LFSGray,
    rounded corners=3pt,
    dashed,
    line width=0.7pt,
    inner sep=6pt
  },
  arrow/.style={
    -{Latex[length=2.4mm,width=1.6mm]},
    draw=LFSGray,
    line width=0.9pt
  },
  data arrow/.style={
    arrow,
    draw=LFSBlue
  },
  control arrow/.style={
    arrow,
    draw=LFSOrange,
    dashed
  },
  residual/.style={
    arrow,
    draw=LFSGreen,
    rounded corners=4pt
  },
  label text/.style={
    font=\scriptsize,
    text=LFSGray,
    fill=white,
    inner sep=1.5pt
  },
  title/.style={
    font=\bfseries,
    text=LFSGray,
    align=center
  }
}


\begin{document}
\begin{tikzpicture}[my picture, node distance=5mm and 6.5mm]

% 4个训练演进阶段：预训练 -> 微调 -> 偏好对齐 -> 可验证推理

% 阶段 1：预训练与规模定律
\node[storage, text width=32mm, minimum height=14mm] (s1_top) {
  \textbf{数据工程与规模定律}\\[0.8mm]
  \scriptsize 语料清洗、去重与质量分类\\
  Chinchilla 算力与参数配比（第31、33章）
};
\node[compute, text width=32mm, minimum height=14mm, below=5mm of s1_top] (s1_bot) {
  \textbf{优化器与训练稳定性}\\[0.8mm]
  \scriptsize AdamW、余弦退火与梯度裁剪\\
  Loss Spike 防护与检查点（第14、34章）
};
\begin{scope}[on background layer]
  \node[group, fit=(s1_top)(s1_bot), label={[font=\footnotesize\bfseries, text=MyGreen]above:一、基础预训练阶段}] (g1) {};
\end{scope}

% 阶段 2：指令微调与高效适配
\node[block, text width=32mm, minimum height=14mm, right=7mm of s1_top] (s2_top) {
  \textbf{监督微调（SFT）}\\[0.8mm]
  \scriptsize 多轮对话指令格式化\\
  问答与推理能力激发（第41章）
};
\node[interface, text width=32mm, minimum height=14mm, below=5mm of s2_top] (s2_bot) {
  \textbf{参数高效微调（PEFT）}\\[0.8mm]
  \scriptsize 低秩自适应（LoRA / DoRA）\\
  量化 QLoRA 与前缀微调（第42、43章）
};
\begin{scope}[on background layer]
  \node[group, fit=(s2_top)(s2_bot), label={[font=\footnotesize\bfseries, text=MyBlue]above:二、后微调与高效适配}] (g2) {};
\end{scope}

% 阶段 3：价值与偏好对齐
\node[compute, text width=32mm, minimum height=14mm, right=7mm of s2_top] (s3_top) {
  \textbf{强化学习对齐（RLHF）}\\[0.8mm]
  \scriptsize 人类偏好成对标注\\
  奖励模型与 PPO 策略更新（第44章）
};
\node[interface, text width=32mm, minimum height=14mm, below=5mm of s3_top] (s3_bot) {
  \textbf{直接偏好优化（DPO）}\\[0.8mm]
  \scriptsize 隐式奖励闭式解构造\\
  无奖励模型高效对齐（第45章）
};
\begin{scope}[on background layer]
  \node[group, fit=(s3_top)(s3_bot), label={[font=\footnotesize\bfseries, text=MyOrange]above:三、价值与偏好对齐}] (g3) {};
\end{scope}

% 阶段 4：慢思考可验证推理
\node[block, text width=32mm, minimum height=14mm, right=7mm of s3_top] (s4_top) {
  \textbf{过程奖励模型（PRM）}\\[0.8mm]
  \scriptsize 步骤级正确性自动检验\\
  蒙特卡洛/蒙特卡洛树搜索（第46章）
};
\node[compute, text width=32mm, minimum height=14mm, below=5mm of s4_top] (s4_bot) {
  \textbf{可验证推理强化训练}\\[0.8mm]
  \scriptsize 强化学习自探索演化\\
  测试期计算规模扩展（第46章）
};
\begin{scope}[on background layer]
  \node[group, fit=(s4_top)(s4_bot), label={[font=\footnotesize\bfseries, text=MyPurple]above:四、可验证推理强化}] (g4) {};
\end{scope}

% 阶段间连线
\draw[data arrow] (s1_top.east) -- (s2_top.west);
\draw[data arrow] (s1_bot.east) -- (s2_bot.west);

\draw[data arrow] (s2_top.east) -- (s3_top.west);
\draw[data arrow] (s2_bot.east) -- (s3_bot.west);

\draw[data arrow] (s3_top.east) -- (s4_top.west);
\draw[data arrow] (s3_bot.east) -- (s4_bot.west);

% 底部总结栏
\node[neutral, text width=146mm, minimum height=7mm, below=6mm of $(s2_bot.south)!0.5!(s3_bot.south)$] (bot_summary) {
  \scriptsize \textbf{第三篇训练主线：} 无监督预训练（海量吸收世界知识） $\longrightarrow$ 监督微调（对齐指令格式） $\longrightarrow$ 偏好对齐（契合人类意图） $\longrightarrow$ 过程强化（激发严密推理）
};

\end{tikzpicture}
\end{document}
